\chapter{Stress Health (Stress Management)}

\section{Definition}
Stress is a state of mental or emotional strain resulting from adverse or demanding circumstances. Chronic stress can have significant negative effects on physical and mental health, contributing to a wide range of medical conditions.

\section{Pathophysiology}
The stress response involves activation of the hypothalamic-pituitary-adrenal (HPA) axis and sympathetic nervous system. Cortisol and catecholamines (epinephrine, norepinephrine) are released, increasing heart rate, blood pressure, and blood glucose. Chronic stress leads to persistent elevation of these hormones, causing immune suppression, inflammation, hypertension, insulin resistance, and mental health disorders.

\section{Etiology}
\begin{itemize}
\item Work-related pressure and job insecurity
\item Financial difficulties
\item Relationship problems (marital conflict, divorce, caregiving)
\item Major life changes (moving, job loss, death of a loved one)
\item Academic pressure
\item Health problems (personal or family illness)
\item Trauma (physical, emotional, or sexual abuse)
\item Social isolation
\item Discrimination or minority stress
\end{itemize}

\section{Indications for Evaluation}
\begin{itemize}
\item Anxiety, irritability, or mood swings
\item Fatigue and low energy
\item Difficulty sleeping (insomnia)
\item Difficulty concentrating or making decisions
\item Headaches, muscle tension, or body aches
\item Changes in appetite (overeating or loss of appetite)
\item Increased illness (stress suppresses immune function)
\item Withdrawal from social activities
\item Increased use of alcohol, tobacco, or other substances
\end{itemize}

\section{Related Symptoms}
\begin{itemize}
\item Anxiety disorders
\item Depression
\item Insomnia
\item Headache (tension-type)
\item Hypertension
\item Irritable bowel syndrome
\item Chronic fatigue
\end{itemize}

\section{Self-Care/Home Management}
\begin{itemize}
\item Practice deep breathing, meditation, or mindfulness daily
\item Exercise regularly (30 minutes most days)
\item Maintain a healthy sleep schedule (7-9 hours)
\item Eat a balanced diet and avoid excessive caffeine and alcohol
\item Set boundaries and learn to say no
\item Connect with supportive friends and family
\item Engage in enjoyable hobbies
\item Limit screen time and news consumption
\end{itemize}

\section{Diagnostic Approach}
\begin{itemize}
\item Clinical evaluation with stress assessment tools (Perceived Stress Scale)
\item Screen for anxiety and depression
\item Assess lifestyle factors: sleep, exercise, diet, substance use
\end{itemize}

\section{Treatment/Management Overview}
\begin{itemize}
\item Cognitive behavioral therapy
\item Mindfulness-based stress reduction
\item Physical activity and exercise programs
\item Stress management classes or workshops
\item Counseling or psychotherapy for underlying issues
\item Medication for co-occurring anxiety or depression
\item Referral to a mental health specialist
\end{itemize}
